\begin{pastexam}[label={px:gw3-111-ds-04}]{關務三等資訊處理 111 年 第四題}
（申論；10 分）

以文本（text）X=“AGTCATTCGATTC”，樣式（pattern）Y=“ATTC”兩字串為例，請問使用暴力比較/窮舉法（exhaustive
search）中的樣式前向法（forward）及後向法（backward）各需比較幾次？（10 分）
\end{pastexam}

\begin{pastsoln}
答案：找出所有出現位置時，前向法比較 $\mathbf{17}$ 次、後向法 $\mathbf{18}$ 次。 （若找到第一個出現位置就停止，
前向法 $9$ 次、後向法 $10$ 次。）

\textbf{解讀}。 暴力法（資結書第~27.1~節的樸素比對）把樣式依序對齊文本的每個位置 $s = 0, 1, \ldots, n - m$，逐字元
比較，遇到不相同的字元就換下一個位置。 「前向」「後向」指每個位置上\textbf{比較字元的方向}：前向法從樣式的第一個
字元往後比，後向法從樣式的最後一個字元往前比。 一次「比較」是拿一個樣式字元和一個文本字元比一次。 題目沒有說找到
後要不要停，一般「字串比對」要找出所有出現的位置，所以主答案比完全部位置；只找第一個的次數也列在後面。

$X$ 長 $n = 13$，$Y$ 長 $m = 4$，位置 $s = 0..9$ 共 $10$ 個。 每個位置對齊的子字串 $X[s..s+3]$ 和兩種方法的比較次數：
\begin{center}
\small
\begin{tabular}{cclclc}
\toprule
$s$ & $X[s..s+3]$ & 前向（依序比 A、T、T、C） & 次數 & 後向（依序比 C、T、T、A） & 次數 \\
\midrule
0 & AGTC & A$=$A，G$\ne$T & 2 & C$=$C，T$=$T，G$\ne$T & 3 \\
1 & GTCA & G$\ne$A & 1 & A$\ne$C & 1 \\
2 & TCAT & T$\ne$A & 1 & T$\ne$C & 1 \\
3 & CATT & C$\ne$A & 1 & T$\ne$C & 1 \\
4 & ATTC & 全部相同（找到） & 4 & 全部相同（找到） & 4 \\
5 & TTCG & T$\ne$A & 1 & G$\ne$C & 1 \\
6 & TCGA & T$\ne$A & 1 & A$\ne$C & 1 \\
7 & CGAT & C$\ne$A & 1 & T$\ne$C & 1 \\
8 & GATT & G$\ne$A & 1 & T$\ne$C & 1 \\
9 & ATTC & 全部相同（找到） & 4 & 全部相同（找到） & 4 \\
\midrule
合計 & & & 17 & & 18 \\
\bottomrule
\end{tabular}
\end{center}
（位置從 $0$ 編號；$X$ = A G T C A T T C G A T T C，樣式出現在 $s = 4$ 和 $s = 9$。）

\begin{itemize}
\item 前向法：$2 + 1 + 1 + 1 + 4 + 1 + 1 + 1 + 1 + 4 = \mathbf{17}$ 次。
\item 後向法：$3 + 1 + 1 + 1 + 4 + 1 + 1 + 1 + 1 + 4 = \mathbf{18}$ 次。
\end{itemize}
兩者只差在 $s = 0$：前向法第二個字元（G 對 T）就失敗，後向法要比到第三個才失敗。 若找到第一個出現位置（$s = 4$）就
停止，前向法是 $2 + 1 + 1 + 1 + 4 = 9$ 次，後向法是 $3 + 1 + 1 + 1 + 4 = 10$ 次。

最壞情況下暴力法每個位置都要比 $m$ 次，共 $(n - m + 1)m$ 次（這裡是 $10 \times 4 = 40$）；實際上大多數位置在第一個
字元就失敗，所以次數遠小於最壞情況。
\end{pastsoln}
